% from ega3-1-fr.tex
\oldpage[0\textsubscript{III}]{71}
Cela étant, supposons \emph{b)} vérifiée~; pour tout $p\in\mathbb Z$, et
tout $k\geq p$, il existe un entier $\mathrm L(p,k)$ tel que le second
membre de (\hyperref[0.13.4.3.2-fr]{13.4.3.2}) soit constant pour
$h\geq\mathrm L(p,k)$~; comme $\operatorname{gr}^p(A_k)=0$ pour $p<k_0$,
il n'y a (pour $k$ donné) qu'un nombre fini de $\mathrm L(p,k)$ non nuls
lorsque $p$ parcourt l'ensemble des entiers $\leq k$. Soit
$\mathrm L(k)=m$ le plus grand de ces entiers~; pour tout $h\geq m$, on a
$A_{kh}\subset A_{km}$, et par définition de $m$, l'injection canonique
$A_{kh}\to A_{km}$ définit un \emph{isomorphisme}
$\operatorname{gr}^\bullet(A_{kh})\xrightarrow{\sim}
\operatorname{gr}^\bullet(A_{km})$~; comme les filtrations sont
\emph{finies}, on en conclut que l'injection précédente est elle-même
bijective (Bourbaki, \emph{Alg. comm.}, chap.~III, §~2,
n\textsuperscript{o}~8, th.~1), ce qui prouve que $\mathbf A$ vérifie
(ML).
\end{proof}

\begin{env}[13.4.4]
\label{0.13.4.4-fr}
Supposons que dans $C$ la limite projective
$A=\varprojlim_k A_k$ existe. Dans les définitions de
(\hyperref[0.13.4.1-fr]{13.4.1}), on peut alors remplacer $A_k$ par $A$,
et la filtration ainsi définie sur $A$ est encore telle que
\begin{equation}
\label{0.13.4.4.1-fr}
\tag{13.4.4.1}
\operatorname{gr}^p(A)=\operatorname{gr}^p(A_p)
\cap\operatorname{Im}(A\longrightarrow A_p).
\end{equation}
\end{env}

\begin{corollary}[13.4.5]
\label{0.13.4.5-fr}
\emph{Supposons que $C$ soit la catégorie des groupes abéliens. Si le
système projectif $\mathbf A$ vérifie (ML) et si
$A=\varprojlim_k A_k$, on a pour tout $p\in\mathbb Z$ un isomorphisme
canonique}
\begin{equation}
\label{0.13.4.5.1-fr}
\tag{13.4.5.1}
\operatorname{gr}^p(A)\xrightarrow{\sim}
\varprojlim_k\operatorname{gr}^p(A_k).
\end{equation}
\end{corollary}
\begin{proof}
En effet, on a
$\operatorname{Im}(A_k\to A_p)=\operatorname{Im}(A\to A_p)$ dès que $k$
est assez grand (Bourbaki, \emph{Top. gén.}, chap.~II,
3\textsuperscript{e} éd., §~3, n\textsuperscript{o}~5, th.~1), et la
conclusion résulte de
(\hyperref[0.13.4.1.3-fr]{13.4.1.3}) et
(\hyperref[0.13.4.4.1-fr]{13.4.4.1}).
\end{proof}

\subsection{Limites projectives de suites spectrales de complexes filtrés}
\label{subsection:0.13.5-fr}
\begin{env}[13.5.1]
\label{0.13.5.1-fr}
Soient $C$ une catégorie abélienne, $X^\bullet$ un complexe d'objets de
$C$ muni d'une \emph{filtration}
$(\mathrm F^p(X^\bullet))_{p\in\mathbb Z}$ telle que
$\mathrm F^{p_0}(X^\bullet)=X^\bullet$ pour un indice $p_0$.
Considérons pour chaque $k\in\mathbb Z$ le complexe
$X_k^\bullet=X^\bullet/\mathrm F^{k+1}(X^\bullet)$~; il est canoniquement
muni de la filtration formée des
$\mathrm F^p(X_k^\bullet)=
\mathrm F^p(X^\bullet)/\mathrm F^{k+1}(X^\bullet)$ pour $p\leq k$ et des
$\mathrm F^p(X_k^\bullet)=0$ pour $p\geq k+1$. En outre, on a des
morphismes canoniques $X_{k+1}^\bullet\to X_k^\bullet$, qui font de
$\mathbf X^\bullet=(X_k^\bullet)_{k\in\mathbb Z}$ un \emph{système
projectif de complexes filtrés d'objets de $C$}. On notera que ce système
projectif est \emph{strict} et tel que $X_k^\bullet=0$ pour $k<p_0$.
\end{env}
\begin{env}[13.5.2]
\label{0.13.5.2-fr}
Considérons plus généralement un système projectif \emph{strict}
$\mathbf X^\bullet=(X_k^\bullet)_{k\in\mathbb Z}$ de complexes d'objets de
$C$, \emph{limité inférieurement}~; considérons sur chaque
$X_k^\bullet$ la filtration définie dans
(\hyperref[0.13.4.1-fr]{13.4.1}) (en se plaçant dans la catégorie abélienne
des complexes de $C$ limités inférieurement). Les
$X_k^\bullet\to X_p^\bullet$ $(p\leq k)$ deviennent des morphismes de
complexes \emph{filtrés}, à filtrations finies. Le caractère fonctoriel des
suites spectrales des complexes filtrés
(\hyperref[0.11.2.3-fr]{11.2.3}) montre que les morphismes de définition du
système projectif $\mathbf X^\bullet$ fournissent des morphismes faisant de
$\mathrm E(\mathbf X^\bullet)=(\mathrm E(X_k^\bullet))$ un \emph{système
projectif de suites spectrales}.
\end{env}
\begin{lemma}[13.5.3]
\label{0.13.5.3-fr}
\emph{Supposons que le système projectif
$\mathbf X^\bullet=(X_k^\bullet)_{k\in\mathbb Z}$ de complexes filtrés soit
obtenu comme dans (\hyperref[0.13.5.2-fr]{13.5.2}). Alors :}
\begin{enumerate}
\item[\emph{a)}]\emph{Pour $r\geq p-p_0$, on a
$\mathrm B_r^{pq}(X_k^\bullet)=\mathrm B_\infty^{pq}(X_k^\bullet)$ pour
tout $k\in\mathbb Z$.}
\item[\emph{b)}]\emph{Pour $k+1\leq p+r$, on a
$\mathrm Z_r^{pq}(X_k^\bullet)=\mathrm Z_\infty^{pq}(X_k^\bullet)$.}
\item[\emph{c)}]\emph{Pour $k+1\geq p+r$, les morphismes
$\mathrm Z_r^{pq}(X_h^\bullet)\to\mathrm Z_r^{pq}(X_k^\bullet)$ et
$\mathrm B_r^{pq}(X_h^\bullet)\to\mathrm B_r^{pq}(X_k^\bullet)$ sont des
isomorphismes pour tout $h\geq k$.}
\end{enumerate}
\end{lemma}
